\section{Design Goals}
\label{sec:design-goals}

GAWorld is intended as a research instrument rather than a claim that a
language model is, by itself, a valid population model.  Its design goals
therefore concern properties that a researcher can inspect in the running
system.  Each property below has a concrete mechanism in the present
implementation and a boundary that prevents the mechanism from being read as
an empirical guarantee.

\subsection{Persistence across simulated time}

An agent should carry consequences from one decision context into later ones,
rather than being reconstructed from a stateless prompt at every tick.
GAWorld combines a mutable state record with file-backed episodic memory,
daily logs and diaries, schedules, and longer-lived relationship and growth
records.  A stateful run restores the day counter and per-agent memories, while
the memory lifecycle can retrieve, consolidate, summarize, and decay stored
episodes.  This makes persistence testable: a researcher can inspect both the
state before and after a tick and the files used in later recall.  Persistence
does not establish psychological fidelity.  Several memories and summaries
are language-model outputs, configuration can disable lifecycle operations,
and the archived cross-phase evaluation discussed later in the paper is
incomplete.

\subsection{Situatedness in a shared world}

Actions should be conditioned on where and when an agent acts and on what is
locally observable.  The simulation clock exposes day, time, and tick index;
the city map represents named locations and connections; and environment,
local-physical, social, policy, and information mechanisms contribute context
to the cognition loop.  Movement updates location before later stages use the
new setting, and social influence is computed over the current network rather
than over an unconstrained cast of characters.  This is situated computation,
not a calibrated digital twin: map detail, physical dynamics, and agent access
to information are abstractions selected by configuration.

\subsection{Composability without concealing control flow}

Researchers should be able to change an experimental mechanism without
rewriting the complete agent loop.  GAWorld provides three bounded extension
surfaces: a configurable stage sequence for cognition, an event bus with
observe/contribute/filter semantics, and a registry for lifecycle-managed
plugins.  The audited built-in set contains nine plugins: interventions,
skills, interests, life events, economy, local physical perception, real work,
dynamic behavior, and spatial preferences.  Composability is not unrestricted
hot swapping: stages exchange a shared dictionary, while memory, environment,
and social-network logic retain direct compatibility paths.

\subsection{Controllability for counterfactual experiments}

A research run should expose explicit points at which a scenario or
intervention is introduced.  Configuration fixes profiles, city assets,
random seed, mechanism flags, schedules, and policy events; the intervention
plugin contributes a recommendation feed through the event bus and records
its domain-specific metrics.  Baseline/event workflows can then align traces
produced under different scenario specifications.  The kernel also defines a
Controller API for validators and named runtime interventions.  Structured
movement requests pass through its validator chain, and the standard runtime
interventions update agent state, update live configuration, or queue agent
removal with audit records.  Other selected actions are not universally
mediated, so configuration, scheduled events, hooks, and experiment scripts
remain important control surfaces.

\subsection{Inspectability from decision to artifact}

A claim about an artificial society should be traceable to artifacts below the
level of a prose summary.  GAWorld writes per-agent logs and memories,
environment timelines, state histories, plugin-owned outputs, and experiment
comparisons; its interfaces expose portions of those records for inspection
and replay.  The kernel Recorder additionally implements a clock-stamped JSONL
stream and records Controller interventions and denials.  That unified stream
is not yet the main action/trace path: the established simulator continues to
write several specialized formats directly.  Inspection therefore requires
an artifact ledger that names the producing path and schema, and replayability
must not be inferred merely from the presence of logs.

\subsection{Backward-compatible modularization}

The system must remain usable while a monolithic research prototype is
decomposed.  A typed \texttt{Agent} view shares, rather than copies, the legacy
dictionary; \texttt{EventBus} preserves the legacy hook signature; and the
12-stage pipeline was carved from the existing per-agent step while retaining
the shared step dictionary.  These adapters allow migrated and legacy modules
to coexist and make migration state observable.  They also retain weakly typed
surfaces and shared mutable state.  Backward compatibility is thus a transition
strategy, not evidence that modularization is complete or that every old
configuration produces an identical trajectory.
